\documentclass[12pt]{book}

\usepackage{amscd}                 %% Celui-ci provoque des erreurs!!!
\usepackage[utf8]{inputenc}
\usepackage[frenchb]{babel}
\usepackage{pictex, color}
\usepackage{amsmath, amssymb, amstext, amsthm}
\usepackage{enumerate}

\usepackage{pstricks,pstricks-add,pst-node,pst-tree}
\usepackage{auto-pst-pdf}

\def\eqdef {\stackrel {\rm def}{=}}   % Egal par definition.
\def\eqas  {\stackrel {\rm p.s.}{=}}  % Egal a.s.
\def\toas  {\stackrel {\rm p.s.}{\to}}   % ---> conv. a.s.
\def\tod   {~\Rightarrow~}               % ---> conv. en distribution.
\def\eff {\mbox{\rm Eff}}
\def\iid {\mbox{\rm i.i.d.{}}}
\def\MSE {\mbox{\rm MSE}}
\def\RE  {\mbox{\rm RE}}
\def\rmis {{\rm is}}

\def\cF {\mathcal{F}}
\def\cH {\mathcal{H}}
\def\cQ {\mathcal{Q}}
\def\cR {\mathcal{R}}
\def\cS {\mathcal{S}}
\def\cX {\mathcal{X}}

\def\Var {\text{Var}}
\def\Cov {\text{Cov}}
\def\To {\Rightarrow}
\def\RR {\mathbb{R}}
\def\bA {\boldsymbol{A}}
\def\bL {\boldsymbol{L}}
\def\bI {\boldsymbol{I}}
\def\bmu {\boldsymbol{\mu}}
\def\bR {\boldsymbol{R}}
\def\bSigma {\boldsymbol{\Sigma}}
\def\bU {\boldsymbol{U}}
\def\bY {\boldsymbol{Y}}
\def\bX {\boldsymbol{X}}
\def\bx {\boldsymbol{x}}
\def\bxi {\boldsymbol{\xi}}
\def\bZ {\boldsymbol{Z}}
\def\bzero {\boldsymbol{0}}

\def\rmu{{\rm u}}

\def\dDUnif   {{\rm UniformeDiscr\`ete}}
\def\dBer     {{\rm Bernoulli}}
\def\dBinomial{{\rm Binomiale}}
\def\dGeo     {{\rm Géométrique}}
\def\dNegBinomial {{\rm BinomialeN\'egative}}
\def\dPoisson {{\rm Poisson}}
\def\dMultinomial {{\rm Multinomiale}}
%
\def\dUnif    {{\rm Uniforme}}
\def\dNormal  {{\rm Normale}}
\def\dLognormal {{\rm Lognormale}}
\def\dExpon   {{\rm Exponentielle}}
\def\dWeibull {{\rm Weibull}}
\def\dGamma   {{\rm Gamma}}
\def\dBeta    {{\rm Beta}}
\def\dStudent {\mbox{\rm Student-t}}
\def\dChiSquare {{\chi^2}}
\def\dFisher  {\mbox{\rm F}}
\def\dErlang  {{\rm Erlang}}
\def\dPareto  {{\rm Pareto}}
\def\dGumbel  {{\rm Gumbel}}
\def\dPearson {\mbox{\rm Pearson}}
\def\dPearsonV {\mbox{\rm Pearson-5}}
\def\dPearsonVI{\mbox{\rm Pearson-6}}
\def\dRayleigh {{\rm Rayleigh}}
\def\dLogistic {{\rm Logistique}}
\def\dCauchy   {{\rm Cauchy}}

\def\EE {\mathbb{E}}
\def\RR {\mathbb{R}}

\newtheorem{thm}{Theorem}
\newtheorem{prop}{Proposition}

\newcounter {exercisecount}[section]
\renewcommand{\theexercisecount}{\thesection.\arabic{exercisecount}}
\newenvironment {exercise}{}{}%\nlni\begingroup\refstepcounter{exercisecount}%
%\bf \theexercisecount\ \rm }{\xxxend}
\newenvironment {solution}{}{}%\nlni\begingroup\refstepcounter{exercisecount}%

\usepackage{verbatim}

\title{	Introduction à la programmation stochastique\\
	Exercices}

\begin{document}

\maketitle

\chapter{Généralités}

\begin{enumerate}
\item
Démontrez que la valeur de la solution stochastique et la valeur espérée de l'information parfaite sont toujours positives.

Considérons le problèmes général
$$
\min_{x \in X} E[z(x,\bxi)].
$$

Rappelons que la valeur de la solution stochastique (VSS) est définie comme.
$$
VSS = EEV - RP
$$
où $EEV$ est $E[EV]$ et
$$
EV = \min_{x \in X} z(x, \overline{\xi}),\ \overline{\xi} = E[\bxi],
$$
i.e. il s'agit de l'espérance de la valeur de la fonction objectif si la décision est prise en optimisant le problème où l'incertitude a été remplacée par l'espérance des variables aléatoires, et
et $RP$, pour \em{recourse problem}, représente le problème stochastique à résoudre, i.e.
$$
RP = \min_{x \in X} E[z(x, \bxi)].
$$

La valeur espérée de l'information parfaite (EVPI: expected value of perfect information) est elle définie par
$$
EVPI = RP - WS,
$$
où WS, pour \em{wait-and-see}, représente le problème où la décision n'est prise qu'une fois l'incertitude levée, i.e.
$$
WS = E \left[ \min_{x \in X} z(x, \bxi) \right].
$$

Exiger que la VSS est positive revient à demande que
$$
EEV \geq RP
$$
ou
$$
E[ \min_{x \in X} z(x, \overline{\xi}) ] \geq  \min_{x \in X} E[z(x, \bxi)]
$$
Supposons par simplicité que les minimums sont bien définis, et notons
$$
x^*_{EV} \in \arg \min_{x \in X} z(x,\overline{\xi})
$$
Alors, comme $x^*_{EV} \in X$,
$$
E[z(x^*_{EV}, \bxi)] \geq \min_{x \in X} E[z(x, \bxi)],
$$
ce qui correspond à ce que nous recherchions.

Avoir $EVPI \geq 0$ revient à
$$
RP \geq WS
$$
ou
$$
\min_{x \in X} E[z(x, \bxi)] \geq E  \left[ \min_{x \in X} z(x, \bxi) \right]
$$
Notons
$$
x^*_{RP} \in \arg \min_{x \in X} E [z(x,\overline{\bxi})].
$$
Alors, $\forall\, \xi$,
$$
z(x^*_{RP}, \xi) \geq \min_{x \in X} z(x, \xi)
$$
et en passant à l'espérance,
$$
E[z(x^*_{RP}, \bxi)] \geq E \left[ \min_{x \in X} z(x, \xi) \right]
$$
donnant le résultat désiré.

\end{enumerate}

\chapter{Problèmes avec recours}

\mbox{}

\mbox{}

\begin{enumerate}
\item
\begin {exercise}
%(Birge et Louveaux, page 87)
Considérons le programme de seconde étape
\[
Q(x, \xi) = \min_y \lbrace y \ |\ \xi y = 1-x, y \geq 0 \rbrace.
\]
On suppose que $\bxi$ suit une distribution triangulaire sur $[0,1]$, à savoir $P[\bxi \leq u] = u^2$.
\begin{enumerate}[(a)]
\item
Peut-on parler de recours fixe? Pourquoi?
\item
Exprimer $K_2(\xi)$ pour tout $\xi$ dans $[0,1]$.
\item
Exprimer $K_2$, $K_2^P$ et $\cQ$.
\end{enumerate}
\end {exercise}

\item
\begin{exercise}
Considérons le programme de deuxième étape défini comme suit:
\begin{align*}
\min_y\ & 2y_1+y_2\\
\mbox{t.q. } & y_1 - y_2 \leq 2-\bxi x_1,\\
& y_2 \leq x_2,\\
& y_1, y_2 \geq 0.
\end{align*}
Trouvez $K_2(\bxi)$ et $K_2$ pour:
\begin{enumerate}
\item
$\bxi \sim U[0,1]$.
\item
$\bxi \sim \mbox{Poisson}(\lambda),\ \lambda > 0$.
\end{enumerate}
Quelles propriétés attendez-vous pour $K_2$?
Le recours est-il simple ou complet? Si la réponse est non, sous quelles conditions aurions-nous un recours relativement complet?
\end{exercise}

Nous devons donc déterminer des ensembles réalisables.

\mbox{}

Remarquons tout d'abord que comme $y_2 \geq 0$ et $y_2 \leq x_2$, nous
devons avoir $x_2 \geq 0$.
Autrement dit
$$
K_1 = \{ x \,|\, x_1 \in \cR,\ x_2 \geq 0 \}.
$$

\mbox{}

Utilisons la contrainte
\[
y_1 - y_2 \leq 2-\bxi x_1.
\]
Comme $y_2 \leq x_2$, nous avons $y_1 \leq 2-\bxi x_1 + x_2$, et tout
$y_1$ satisfaisant cette inégalité conviendra.
La seule autre contrainte sur $y_1$ étant $y_1 \geq 0$, le problème est réalisable si et seulement si $2-\bxi x_1 + x_2 \geq 0$.

Il suit que
\begin{align*}
K_2(\xi) &= \lbrace x = (x_1, x_2)^T \,|\, x_2 \geq 0,\ \xi x_1 \leq
2+x_2 \rbrace \\
&= \lbrace  x = (x_1, x_2)^T \,|\, x_2 \geq \max(0, \xi x_1 -2 ) \rbrace.
\end{align*}

Pour les distributions considérées, $E[\xi]$ et $E[\xi^2]$ existent et
sont finis. Il suit que
\[
K_2 = K_2^P = \cap_{\xi \in \Xi} K_2(\xi).
\]

\mbox{}

Si $\bxi \sim U[0,1]$, $\xi \leq 1$, et $\xi x_1 \leq 2 + x_2$ si $x_1
\leq 2 + x_2$. Autrement dit, $K_2(\xi) \subseteq K_2(1)$, et
\[
K_2 = \lbrace x \,|\, x_2 \geq 0, x_1 \leq 2 + x_2 \rbrace.
\]

\mbox{}

Si $\bxi \sim \mbox{Poisson}(\lambda)$, $\xi = 0, 1,\ldots,$, de sorte
que $\xi x_1$ n'est pas borné supérieurement, sauf si $x_1 \leq
0$. Dès lors,
\[
K_2 = \lbrace x \,|\, x_1 \leq 0, x_2 \geq 0 \rbrace.
\]

Interprétation. Supposons que toute décision de première étape doit satisfaire
$x_1$, $x_2 \geq 0$ (c'est le cas si par exemple nous travaillons avec la forme standard d'un problème de progammation linéaire en première étape), par conséquent, si $\bxi \sim U[0,1]$,
\[
K_1 \cap K_2 \subseteq \lbrace x \,|\, x_1, x_2 \geq -, x_1 \leq 2+x_2
\rbrace,
\]
et si $\bxi \sim \mbox{Poisson}(\lambda)$,
\[
K_1 \cap K_2 \subseteq \lbrace x \,|\, x_1 = 0, x_2 \geq 0 \rbrace.
\]
Nous pouvons espérer un recours relativement complet dans le premier
cas (cela dépend de $K_1$, qui est non connu), mais le deuxième se présente assez mal!

Le recours n'est pas complet. En effet, si nous choisissons $x_2 < 0$, le problème de second étape n'est pas réalisable.
Cela implique également que le recours n'est pas simple.

\item
Considérons l'approche par pénalisation de la longueur de pas introduite par A. Ruszczynski.
\begin{enumerate}
	\item 
	Pourquoi introduire une telle pénalisation?
	\item 
	Quels en sont les avantages et inconvénients?
	\item 
	Quelle est la différence entre un pas sérieux exact et un pas sérieux approximatif?
\end{enumerate}

Une telle pénalisation a pour but d'éviter que la résolution du problème maître ne produise une solution de première étape ne tenant pas compte d'un point de départ construit à l'aide d'informations sur le problème à résoudre.
Plus précisément, la pénalisation sert à éviter de produire un point trop éloigné de l'estimation disponible à l'itération courante.

Son principal inconvénient est l'introduction d'un terme quadratique dans l'objectif, rendant le coût de résolution du problème maître plus important que dans l'algorithme L-Shaped classique.
Le coût de résolution des problèmes de seconde étape reste par contre inchangé, et le coût total de résolution du problème est généralement moindre, en particulier lorsque le nombre de scénarios est important, grâce aux gains en termes d'itérations.
La diminution du nombre d'itérations vient d'une propension à construire plus raidement des coupes d'optimalité permettant d'isoler la solution optimale.

Un pas sérieux est obtenu quand la solution du problème maître donne une valeur de la fonction objectif du problème à deux étapes moindre que la valeur obtenue avec l'estimation actuelle de la solution.
Le pas est dit approximatif si des coupes d'optimalité ont été produites au cours de l'itération actuelle (aussi la solution du problème maître ne peut être optimale) et sérieux si aucune coupe d'optimalité n'a été ajoutée au cours de l'itération, de sorte que la solution du problème maître est potentiellement optimale.

\item 

\begin{exercise}
%(Birge et Louveaux, page 102)
Soit le programme de seconde étape défini par
\begin{align*}
\min_y\ & 2y_1+y_2\\
\mbox{t.q. } & y_1 + y_2 \geq 1-x_1,\\
& y_1 \geq \bxi - x_1-x_2,\\
& y_1, y_2 \geq 0.
\end{align*}
Montrez que ce programme a un recours complet si $\bxi$ est
d'espérance finie.
\end{exercise}

Le recours est-il complet? Cela revient à exiger que quel que soit le terme de droite, il est possible de trouver une solution réalisable, i.e. $\forall\, z_1, z_2 \in \cR$, $\exists\, y_1, y_2 \geq 0$ tels que
\begin{align*}
y_1 + y_2 &\geq z_1,\\
y_1 &\geq z_2
\end{align*}
Comme il s'agit d'inégalités, nous ne pouvons directement travailler avec des combinaisons linéaires des contraintes.
Toutefois, remarquons que $y_1 = \max\{0, z_2\}$ rend la deuxième contrainte satisfaite, et pour satisfaire la première, il suffit alors de choisir $y_2 = \max\{0, z_1 - y_1\}$. De plus, $y_1$ et $y_2$ sont alors non-négatifs. Le recours est donc complet.

Si nous suivons la définition standard, il faut que pour tout $z \in \cR^{2}$, on puisse
trouver $y \geq 0$ t.q. $Wy = z$. Le recours complet est une propriété de $W$, mais vaut
$W$? De manière équivalente, on veut pos$W = \cR^2$.
Il existe une petite difficulté avec cette approche: le problème n'est pas sous forme standard.
L'idée est en fait de montrer pour pour tout $z \in \cR^2$, il existe
$y \geq 0$ t.q $Wy \geq z$. L'équivalence avec pos$W$ ne fonctionne
plus. Sauf si nous nous ramenons à la forme standard!

\mbox{}

Sous forme standard, le problème devient
\begin{align*}
\min_y\ & 2y_1+y_2\\
\mbox{t.q. } & y_1 + y_2 - s_1 = 1-x_1,\\
& y_1 - s_2 = \bxi - x_1-x_2,\\
& y_1, y_2 \geq 0,\\
& s_1, s_2 \geq 0,
\end{align*}
et il y a recours complet si pour tout $z \in R^2$, il existe $y$, $s
\geq 0$ vérifiant
\[
Wy = z.
\]

Explicitons $Wy = z$:
\[
\begin{pmatrix}
1 & 1 & -1 & 0 \\
1 & 0 & 0  & -1
\end{pmatrix}
\begin{pmatrix} y_1 \\ y_2 \\ s_1 \\ s_2 \end{pmatrix} =
\begin{pmatrix} z_1 \\ z_2 \end{pmatrix}.
\]
Autrement dit, nous devons trouver $y_1$, $y_2$, $s_1$, $s_2 \geq 0$
tel que
\begin{align*}
y_1 + y_2 - s_1 &= z_1. \\
y_1 - s_2 &= z_2.
\end{align*}

La deuxième équation nous donne la possibilité de prendre
\[
\begin{cases}
y_1 = z_2,\ s_2 = 0 & \mbox{si } z_2 \geq 0,\\
y_1 = 0,\ s_2 = -z_2 & \mbox{sinon}.\\
\end{cases}
\]

Nous avons aussi
\[
y_2 - s_1 = z_1 - y_1.
\]
En utilisant nos précédents choix, nous pouvons prendre pour $y_2$ ce
qui suit. Si $z_2 \geq 0$, $y_2 - s_1 = z_1 - z_2$, et nous prenons
\[
\begin{cases}
y_2 = z_1-z_2,\ s_1 = 0 & \mbox{si } z_1-z_2 \geq 0,\\
y_2 = 0,\ s_2 = z_2-z_1 & \mbox{sinon}.\\
\end{cases}
\]
De même, si $z_2 \leq 0$, $y_2 - s_1 = z_1$, et il suffit de choisir
\[
\begin{cases}
y_2 = z_1,\ s_1 = 0 & \mbox{si } z_1 \geq 0,\\
y_2 = 0,\ s_2 = -z_1 & \mbox{sinon}.\\
\end{cases}
\]

Par conséquent, quel que soit $z$, on peut trouver une solution
non-négative au système
\[
W\begin{pmatrix} y \\ s \end{pmatrix} = z,
\]
i.e. le recours est complet!

Prendre $\bxi$ d'espérance finie permet de garantir que $K_2 = \cR^n$.

\item

\begin{exercise}
	Une compagnie doit décider de la commande d'une certaine quantité $x$ d'un produit pour satisfaire une demande $d$, inconnue au moment de la commande.
	Le coût de commande est de $c > 0$ par unité.
	Si la demande $d$ est plus grande que $x$, l'entreprise suffit une pénalité $b \geq 0$ par unité de demande non satisfaite.
	Sinon, si $d < x$, un coût d'inventaire $h$ est encouru par unité excédentaire.
	\begin{enumerate}
		\item
		Formuler le problème comme un problème stochastique à deux étapes.
		\item
		Le recours est-il simple, complet, relativement complet?
		\item
		Que peut-on attendre sur la solution de première étape si $b < c$?
	\end{enumerate}
\end{exercise}

Décision de première étape: $x$. Première étape: commande. Objectif: minimiser les coûts:
\begin{align*}
\min_x \ & c x \\
\mbox{t.q. } & x \geq 0.
\end{align*}
Deuxième étape: satisfaction de la demande. Le coût à minimiser est le coût de pénalité et de le coût d'inventaire:
\begin{align*}
Q(x,d) &= \min_{y^+, y^-} by^+ + hy^- \\
\mbox{t.q. } & y^+ - y^- = d-x\\
& y^+ \geq 0, y^- \geq 0.
\end{align*}
Le problème à deux étapes est donc
\begin{align*}
\min_x \ & c x + E_d [Q(x,d)] \\
\mbox{t.q. } & x \geq 0.
\end{align*}
Le recours est simple, donc également complet et relativement complet.
Si $b < c$, le coût d'achat en première étape est plus important que le coût de pénalité, et la décision optimale de première étape est $x^* = 0$.

\item
%Birge et Louveaux, p 142
On considère le problème suivant:
\begin{align*}
\min_x \ & x_1 + 4 x_2 + E_{\bxi}[Q(x,\xi)] \\
\mbox{t.q. } & x_1 + x_2 = 1,\\
& x \geq 0,
\end{align*}
avec comme problème de recours
\begin{align*}
Q(x,\xi) = \min_y \ & y_1 + 10 y_2 + 10 y_3 \\
\mbox{t.q. } & y_1 + y_2 - y_3 = \bxi + x_1 - 2 x_2,\\
& y_1 \leq 2,\\
&  y_1, y_2, y_3 \geq 0,
\end{align*}
où $x = \begin{pmatrix} x_1 & x_2 \end{pmatrix}^T$, $y
= \begin{pmatrix} y_1 & y_2 & y_3 \end{pmatrix}^T$, $\bxi \sim U[1,3]$.
\begin{enumerate}[(a)]
	\item
	Le recours est-il fixe, simple, complet, relativement complet? Répondez vrai ou faux pour chacun des quatre types de recours, en justifiant brièvement.
	\item
	Pour $x$ et $\xi$ donnés, montrez en discutant sur la valeur de $\xi + x_1 - 2 x_2$ que
	\begin{align*}
	& y^*(x, \xi) = \\
	& \begin{cases}
	y_1 = \xi + x_1 - 2x_2,\ y_2 = y_3 = 0 & \mbox{si } 0 \leq \xi + x_1 -2 x_2
	\leq 2,\\
	y_1 = 2,\ y_2 = \xi + x_1 - 2 x_2 - 2,\ y_3 = 0 & \mbox{si } \xi + x_1 -2 x_2 >
	2,\\
	y_3 = 2x_2 - \xi -x_1,\ y_1 = y_2 = 0 & \mbox{si } \xi + x_1 - 2x_2 < 0.
	\end{cases}
	\end{align*}
	Note: on pourra par exemple discuter en termes de points extrêmes pour étaplir la forme de $y^*(x, \xi)$.
	\item
	Quelle propriété attendue sur $Q(x,\xi)$ est vérifiée en observant la
	valeur $y^*(x,\xi)$?
	\item
	$\mathcal{Q}(x) \overset{\mbox{déf}}{=} E_\xi[Q(x, \xi)]$ est-elle différentiable?
	Pourquoi?
	\item
	Exprimez $K_1$, $K_2(\xi)$, $K_2^P$, $K_2$.
	\item
	Nous considérons maintenant une version discrétisée du problème précédent. Donnez le problème déterministe équivalent (forme étendue) en supposant que l'on discrétise $\bxi$ aux valeurs 1, 2, 3.
\end{enumerate}

\begin{enumerate}[(a)]
\item
Considérons le problème de deuxième étape sous forme standard:
\begin{align*}
Q(x,\xi) &= \min_y q^T(\xi)y \\
\mbox{t.q. } & Wy = h(\xi)-T(\xi)x \\
& y \geq 0.
\end{align*}
Pour obtenir la forme standard, nous ajoutons une variable d'écart, $y_4$, à la deuxième contrainte, pour obtenir
\begin{align*}
Q(x,\xi) = \min_y \ & y_1 + 10 y_2 + 10 y_3 \\
\mbox{t.q. } & y_1 + y_2 - y_3 = \bxi + x_1 - 2 x_2,\\
& y_1 + y_4 = 2,\\
&  y_1, y_2, y_3, y_4 \geq 0.
\end{align*}
Dès lors, on peut écrire
$$
W = \begin{pmatrix} 1 & 1 & -1 & 0 \\ 1 & 0 & 0 & 1 \end{pmatrix},
$$
et donc $W$ ne dépend pas de $\xi$. Autrement dit, le recours est fixe.

On ne peut par contre écrire $W$ sous la forme $\begin{pmatrix} I & -I \end{pmatrix}$, où $I$ est la matrice identité. Le recours n'est donc pas simple.

Le recours est complet si pour tout $z \in \RR^2$, $\exists y \geq 0$ tel que $Wy = z$.
Réécrivons le système $Wy = z$:
\begin{align*}
y_1 + y_2 - y_3 &= z_1 \\
y_1 + y_4 &= z_2.
\end{align*}
Dès que $z_2 < 0$, nous ne pouvons trouvons $y_1, y_4 \geq 0$ satisfaisant la seconde contrainte.
Le recours n'est donc pas complet.

Supposons par contre $(x_1, x_2) = \in K_1$ (non explicitement connu). Nous avons comme solution réalisable, pour $\xi$ donné,
\begin{align*}
y_1 &= 0, \\
y_2 &= (\xi + x_1 - 2 x_2)^+, \\
y_3 &= (-\xi - x_1 + 2 x_2)^+, \\
y_4 & = 2,
\end{align*}
où $x^+ = \max \{ x, 0 \}$. Dès lors, le recours est relativement complet.

\item

Considérons la valeur de $\xi + x_1 - 2 x_2$.
Si $\xi + x_1 - 2 x_2 < 0$, la seule solution extrême réalisable est 
$$
y_1 = y_2 = 0, \ y_3 = -\xi - x_1 + 2 x_2, \ y_4 = 2
$$
et est donc optimale.
Si $\xi + x_1 - 2 x_2 \geq 0$, nous devrons avoir pour un point extrême $y_3 = 0$, et nous avons
$$
y_1 + y_2 = \xi + x_1 - 2 x_2,\ y_1 \leq 2.
$$
Comme le coefficient de coût de $y_1$ est 1, et celui de $y_2$ est 10, il suit que $y_1$ devrait être choisi comme
$$
y_1 = \min \{ \xi + x_1 - 2 x_2, 2 \}.
$$
Nous avons donc deux cas à discuter: $\xi + x_1 - 2 x_2 \leq 2$ et $\xi + x_1 - 2 x_2 \geq 2$.
Si $\xi + x_1 - 2 x_2 \leq 2$, nous aurons
$$
y_1 = \xi + x_1 - 2 x_2,\ y_2 = 0,
$$
tandis $\xi + x_1 - 2 x_2 \geq 2$,
$$
y_1 = 2,\ y_2 = \xi + x_1 - 2 x_2 - 2,
$$
comme escompté.

\item

La fonction $Q(x, \xi)$ est convexe et linéaire par morceaux par rapport à la décision de première étape $x$.

\item 

Oui, car $\bxi$ est une variable aléatoire continue ayant un moment d'ordre 2 fini.

\item

De l'énoncé, nous avons directement
$$
K_1 = \{ x \,|\, x_1+x_2 = 1, x_1 \geq 0, x_2 \geq 0 \}.
$$

Du point (b), on voit que quel soit $\xi$, $K_2(\xi) = \RR^2$. Par conséquent, $K_2^P = \cap_{\xi \in \Xi} K_2(\xi) = \RR^2$.

\item

L'équivalent déterministe peut se formuler comme
\begin{align*}
\min_x \ & x_1 + 4 x_2 + \sum_{i=1}^3 p_i (y_{1i}+10y_{2i}+10y_{3i}) \\
\mbox{t.q. } & x_1 + x_2 = 1,\\
& y_{1i} + y_{2i} - y_{3i} = \xi_i + x_1 - 2 x_2, \quad i = 1,2,3\\
& y_{1i} \leq 2,\\
& x, y_{1i}, y_{2i}, y_{3i} \geq 0, \quad i = 1,2,3,
\end{align*}
où $i$ désigne l'indice de scénario.
Nous supposons les scénarios équiprobables, aussi $p_i = 1/3$ pour tout $i$ et $\xi_1 = 2$, $\xi_2 = 2$, $\xi_3 = 3$.

\item

Suppose that we have to make a decision about the amount of a good $X$ to
produce.
Each unit of $X$ that we make costs us \$2, and we have to meet demand from customers in the next time period.
However the demand is stochastic, with a discrete probability distribution. There are $S$ possible demand values, denoted by $D_s$, $s = 1,\ldots,S$, with the respective probabilities $p_s$, $s = 1,\ldots,S$.

We have the flexibility to buy in the product from an external supplier to meet
observed customer demand but this costs us \$3 per unit.
How much should we choose to make now before we know what customer
demand is?
Formulate the problem as a two-stage stochastic program, minimizing the expected production costs.

Let $x \geq 0$ be the number of units of $X$ to produce now (at the first stage)
Let $y_s \geq 0$ be the number of units of $X$ to buy from the external supplier at the
second stage in scenario $s$ when the stochastic realization of the demand is $D_s$.
Then the constraints to ensure that demand is always satisfied are
$$
x + y_s \geq D_s,\ s = 1,\ldots,S,
$$
and the stochastic program is
\begin{align*}
\min\ & 2x + \sum_{s = 1}^S p_s 3 y_s \\
\mbox{s.t. } &  x + ys \geq D_s,\ s = 1,\ldots,S \\
& x \geq 0 \\
& y_s \geq 0, s = 1,\ldots,S.
\end{align*}

\end{enumerate}

\end{enumerate}

\chapter{Problèmes multi-étapes}

\begin{enumerate}
\item
\begin{exercise}
	Considérons le programme stochastique multi-étapes stylisé suivant (Heitsch et al., 2006),
	cherchant à déterminer une stratégie optimale d'achat sous des coûts incertains $\bxi_t$,
	sur les périodes $t = 1,\ldots,T$. Les décisions de montants à acheter sont donnés par $x_t$,
	et nous souhaitons minimiser les coûts espérés de manière à atteindre un montant prescrit
	$a$ à l'horizon $T$:
	\begin{align*}
	\min_{\bx_t}\ & E\left[ \sum_{t=0}^T \bxi_t \bx_t \right],\\
	\mbox{s.t. } & s_t-s_{t-1} = \bx_t,\ \forall t = 0,\ldots,T,\\
	& s_{-1} = 0,\ s_T = a,\ \bx_t \geq 0,\ s_t \geq 0,
	\end{align*}
	où $s_t$ est une variable d'état  contenant le montant au temps $t$.
	
	Nous considérons 3 étapes ($t = 0, 1, 2$), un arbre binaire, et $\xi_0 = 10$.
	La valeur monte de 1 avec une probabilité 0.6, et descent de 1 avec une probabilité 0.4.
	
	\begin{enumerate}
		\item 
		Représentez l'arbre de scénarios.
		\item
		Écrivez le problème
		\begin{enumerate}
			\item
			sous forme étendue
			\item
			à l'aide d'une décomposition par scénarios.
		\end{enumerate}
	\end{enumerate}
	
	%\begin{thebibliography*}{1}
	
	%	\bibitem{HeitRomiStru06}
	{\small Référence:
		H.~Heitsch and W.~Römisch and C.~Strugarek, Stability of multistage stochastic
		programs, {\em SIAM Journal on Optimization}, 17(2):511--525 (2006).
	}
	
	%\end{thebibliography*}
\end{exercise}

\begin{center}
%\begin{footnotesize}
%\scalebox{0.85}
{
\pstree[treemode=R,levelsep=22ex]{\Tcircle{$\xi_0 = 10$}}{
\pstree{\Tcircle{$\xi_1 = 9$}
\ncput{0.4}}
{\Tcircle{$\xi_2 = 8$}
\trput{0.4}
%~[tnpos=r]{\quad\bf{\large {\blue $1$}}}
%					[tnpos=b]{\quad\bf{\large {\blue $1$}}}
\Tcircle{$\xi_2 = 10$}
\trput{0.6}
%~[tnpos=r]{\quad\bf{\large {\blue $2$}}}
}
\pstree{\Tcircle{$\xi_1 = 11$}\ncput{0.6}}
{
\Tcircle{$\xi_2 = 10$}%~[tnpos=r]{\quad\bf{\large {\blue $3$}}}
\trput{0.4}
\Tcircle{$\xi_2 = 12$}%~[tnpos=r]{\quad\bf{\large {\blue $4$}}
\trput{0.6}}
}
}
	%\end{footnotesize}
\end{center}

\begin{enumerate}[a)]
\item
Forme étendue.

Nous attachons une décision à chaque noeud: $x_{0}$, $x_{1,1}$, $x_{1,2}$, $x_{2,1}$, $x_{2,2}$, $x_{2,3}$, $x_{2,4}$.

\begin{align*}
\min_{\bx_t}\ & 10x_0+0.4(9x_{1,1}+0.4*8x_{2,1}+0.6*10x_{2,2}) \\
& +0.6(11x_{1,2}+0.4*10x_{2,3}+0.6*12x_{2,4}),\\
\mbox{s.t. } & s_0 = x_0, \\
& s_{1,1} = x_{1,1}+s_0, s_{1,2} = x_{1,2}+s_0, \\
& s_{2,1} = x_{2,1}+s_{1,1}, s_{2,2} = x_{2,2}+s_{1,1}, \\
& s_{2,3} = x_{2,3}+s_{1,2}, s_{2,4} = x_{2,4}+s_{1,2}, \\
& s_{-1} = {0},\ s_{2,1} = s_{2,2} = s_{2,3} = s_{2,4} = a,\\
& x_0, x_{1,1}, x_{1,2}, x_{2,1}, x_{2,2}, x_{2,3}, x_{2,4} \geq 0,\ s_t \geq 0.
\end{align*}
\item
Décomposition par scénarios.

Il y a 4 scénarios: $x_{0,1}$, $x_{0,2}$, $x_{0,3}$, $x_{0,4}$,
$x_{1,1}$, $x_{1,2}$, $x_{1,3}$, $x_{1,4}$,
$x_{2,1}$, $x_{2,2}$, $x_{2,3}$, $x_{2,4}$.

\begin{align*}
\min_{\bx_t}\ &
0.16(10x_{0,1} + 9x_{1,1} + 8x_{2,1})
+ 0.24(10x_{0,2} + 9x_{1,2} + 10x_{2,2}) \\
& + 0.24(10x_{0,3} + 11x_{1,3} + 10x_{2,3})
+0.36(10x_{0,4} + 11x_{1,4} + 12x_{2,4}) \\
\mbox{s.t. } & x_{0,1} = x_{0,2} = x_{0,3} = x_{0,4}, \\
& s_{1,1} - x_{0,1} = x_{1,1},
s_{1,2} - x_{0,2} = x_{1,2},
x_{1,1} = x_{1,2} \\
& s_{1,3} - x_{0,3} = x_{1,3},
s_{1,4} - x_{0,4} = x_{1,4},
x_{1,3} = x_{1,4} \\
& s_{1,1} - s_0 = x_{1,1} = x_{1,2},
s_{1,2} - s_0 = x_{1,3} = x_{1,4}, \\
& s_{2,1} = x_{2,1}+s_{1,1}, s_{2,2} = x_{2,2}+s_{1,1},
 s_{2,3} = x_{2,3}+s_{1,2}, s_{2,4} = x_{2,4}+s_{1,2}, \\
& s_{-1} = {0},\ s_{2,1} = s_{2,2} = s_{2,3} = s_{2,4} = a,\\
& x_0, x_{1,1}, x_{1,2}, x_{2,1}, x_{2,2}, x_{2,3}, x_{2,4} \geq 0,\ s_t \geq 0.
\end{align*}
\end{enumerate}

\item

Dans la méthode de régularisation par pénalisation proposée by Ruszczy\'nski, quelle est la différence entre pas sérieux exact et pas sérieux approximatif? Pourquoi doit-on les distinguer?

Très brièvement, la différence principale est si le pas a le potentiel de fournir la solution optimale ou non. Dans les deux cas, nous acceptons un pas si nous obtenons une meilleure solution que celle connue actuellement, mais avec le pas approximatif, nous savons que la solution n'est pas optimale, alors qu'avec le pas exact, elle pourrait être optimale.

De manière plus détaillée, en considérant la méthode régularisée:
\begin{itemize}
\item 
À l'étape 4, nous vérifions si aucune coupe d'optimalité n'a été ajoutée. Si tel est la cas, dans l'algorithme standard, cela impliquerait que la solution du problème principal, résolue à l'étape 1, est optimale, mais ici, en raison de la pénalité, nous ne savons pas encore si c'est le cas.
Nous prenons néanmoins la solution actuelle de première étape comme nouvelle estimation de la décision de première étape et nous revenons à l’étape 1. Le dernier test de cette étape vérifie si la résolution du problème principal donne que $x^\nu$ est égal à $a^{\nu}$, de sorte que la pénalité disparait et comme nous n’avons pas eu de nouvelles coupes d’optimalité, la solution est optimale. Sinon, nous observons que le pénalité a empêché d’obtenir la solution optimale et nous continuons.
\item
L’étape 5 ne peut être atteinte que si la condition de l’étape 4 est violée, ou en d’autres termes, si nous avons ajouté au moins une nouvelle coupe d'optimalité à l’étape 3. Si nous avons ajouté de nouvelles coupes, la solution ne peut pas être optimale.
Elle pourrait cependant fournir un meilleur coût lors de l’évaluation de la fonction objectif complète à deux étapes. Si tel est le cas, nous déclarons que notre solution de première étape est meilleure car elle fournit un coût objectif plus bas et nous la prenons comme nouvelle estimation $a^{\nu}$ de la solution, mais nous savons déjà qu'elle n’est pas optimal et nous devons continuer.
\end{itemize}

\item
% http://uclengiechair.be/wp-content/uploads/2017/05/LShaped.pdf
Considérez le problème
\begin{align*}
\min\ & 100x_1 + 150x_2 + \EE_{\bxi} [q_1y_1 + q_2y_2] \\
\mbox{t.q. }
& x_1 + x_2 \leq 120 \\
& 6y_1 + 10y_2 \leq 60x_1 \\
& 8y_1 + 5y_2 \leq 80x_2 \\
& y_1 \leq d_1,\ y_2 \leq d_2 \\
& x_1 \geq 40,\ x_2 \geq 20, \\
& y_1,\ y_2 \geq 0
\end{align*}
sous deux scénarios $\bxi = (d_1, d_2, q_1, q_2)$.
Soit $\xi_1 = (500, 100, -24, -28)$, avec la probabilité $p_1 = 0.4$, et $\xi_2 = (300, 300, -28, -32)$, avec la probabilité $p_2 = 0.6$.

Éxécutez la première itération de la méthode $L$-shaped, simple coupe, en montrant que les solutions associées aux deux scénarios sont, à la première itération, respectivement $y^* = (137.5, 100)$ et $y^* = (80, 192)$.
Pour déterminer les solutions duales, on pourra se baser sur les écarts de complémentarité, indiquant qu'une variable duale ne peut être non nulle que si la contrainte primale associée est active.

Le problème maître à résoudre à la première itération du L-Shaped est
\begin{align*}
\min\ & 100x_1 + 150x_2  \\
\mbox{t.q. }
& x_1 + x_2 \leq 120 \\
& x_1 \geq 40,\ x_2 \geq 20.
\end{align*}
On peut immédiatement voir que la solution optimale est $x_1^* = 40$, $x_2^* = 20$.

Le problème de deuxième étape, étant donné $\xi$, est
\begin{align*}
\min_y\ & q_1y_1 + q_2y_2 \\
\mbox{t.q. }
& 6y_1 + 10y_2 \leq 60x_1 \quad (\pi_1) \\
& 8y_1 + 5y_2 \leq 80x_2 \quad (\pi_2) \\
& y_1 \leq d_1  \quad (\pi_3) \\
& y_2 \leq d_2 \quad (\pi_4) \\
& y_1,\ y_2 \geq 0
\end{align*}
où $\pi$ représente le vecteur des variables duales.
Son dual est
\begin{align*}
\max_{\pi} \ &2400\pi_1 + 1200\pi_2 + d_1\pi_3 + d_2 \pi_4 \\
\mbox{t.q. }
& 6\pi_1 + 8\pi_2 +\pi_3 \leq q_1 \\
& 10\pi_1 + 5\pi_2+\pi_4 \leq q_2 \\
& \pi \leq 0
\end{align*}

Comme $x_1$, $x_2 \geq 0$, ce programme est toujours réalisable, autrement dit le recours est relativement complet.
Nous n'avons donc pas à considérer la génération de coupe d'optimalité.

Résolvons chacun des deux problèmes de deuxième étape, étant donné la solution de première étape $x_1^* = 40$, $x_2^* = 20$.

Premier scénario:
\begin{align*}
\min_y\ & -24y_1 -28y_2 \\
\mbox{t.q. }
& 6y_1 + 10y_2 \leq 2400 \\
& 8y_1 + 5y_2 \leq 1600 \\
& y_1 \leq 500,\ y_2 \leq 100 \\
& y_1,\ y_2 \geq 0
\end{align*}

Nous devons avoir $y_1 < 500$ comme sinon, la deuxième contrainte est violée.

Supposons $y_2 = 100$. Alors
\begin{align*}
y_1 &= \min \left\{ \frac{2400-1000}{6}, \frac{1600-500}{8} \right\} \\
&= \left\{ \frac{700}{6}, \frac{1100}{8} \right\} \\
&= \{ 115+10/6, 137.5 \} \\
&=  137.5
\end{align*}

Considérons le cas $y_1 < 500$, $y_2 < 100$. Ce sera le cas si la solution consiste à rendre les deux premières contraintes actives. Dans cas,
\begin{align*}
3y_1 + 5y_2 = 1200 \\
8y_1 + 5y_2 = 1600
\end{align*}
Dans cas, $y_1 = 80$, $y_2 = 192$, ce qui viole la contrainte $y_2 = 100$.

La solution optimale est donc $y_1^* = 137.5$, $y_2^* = 100$.

On en déduit directement des écarts de complémentarité que $\pi_1^* = \pi_3^* = 0$, et $\pi_2* = -24/8 = -3$, $\pi_4^* = -28-15=-13$.

Second scénario:
\begin{align*}
\min_y\ & -28y_1 -32y_2 \\
\mbox{t.q. }
& 6y_1 + 10y_2 \leq 2400 \\
& 8y_1 + 5y_2 \leq 1600 \\
& y_1 \leq 300,\ y_2 \leq 300 \\
& y_1,\ y_2 \geq 0
\end{align*}

Prendre $y_1 = 300$ ou $y_2 = 300$ mène à la violation d'une des deux premières contraintes.
La solution s'obtient en résolvant le système
\begin{align*}
3y_1 + 5y_2 = 1200 \\
8y_1 + 5y_2 = 1600
\end{align*}
dont la solution a déjà été obtenue: $y_1^* = 80$, $y_2^* = 192$. Dès lors, $\pi_3* = \pi_4^* = 0$, et
\begin{align*}
6\pi_1 + 8\pi_2 &= -28 \\
10\pi_1 + 5\pi_2 &= -32
\end{align*}
ou
\begin{align*}
3\pi_1 + 4\pi_2 &= -14 \\
10\pi_1 + 5\pi_2 &= -32
\end{align*}
ou encore
\begin{align*}
15\pi_1 + 20\pi_2 &= -70 \\
40\pi_1 + 20\pi_2 &= -128
\end{align*}
On en tire $25\pi_1^* = -58$ ou encore $\pi_1^* = -2.32$ et $5\pi_2^* = -32+10\times 58/25 = -32+116/5 = -$,
ce qui donne $\pi_2^* = -1.76$.

Afin de construire la coupe d'optimalité, reprenons les contraintes de deuxième étape sous la forme
$$
Wy \leq h(\xi)-T(\xi)x
$$
Nous avons
$$
Wy \leq h(\xi_1) - T(\xi_1)x
$$
avec
$$
h(\xi_1) = \begin{pmatrix} 0 \\ 0 \\ 500 \\ 100 \end{pmatrix} \quad 
T(\xi_1) = \begin{pmatrix} -60 & 0 \\ 0 & -80 \\ 0 & 0 \\ 0 & 0  \end{pmatrix}
$$
et
$$
h(\xi_2) = \begin{pmatrix} 0 \\ 0 \\ 300 \\ 300 \end{pmatrix} \quad 
T(\xi_2) = \begin{pmatrix} -60 & 0 \\ 0 & -80 \\ 0 & 0 \\ 0 & 0  \end{pmatrix}
$$
Nous remarquons que $T(\xi_1) = T(\xi_2)$, et notons cette matrice $T$.
Dès lors
\begin{align*}
e_1 &= 0.4\pi_1^Th(\xi_1) + 0.6\pi_2^Th(\xi_2) \\
&= 0.4  \begin{pmatrix} 0 & -3 & 0 -13 \end{pmatrix} \begin{pmatrix} 0 \\ 0 \\ 500 \\ 100 \end{pmatrix} + 
0.6  \begin{pmatrix} -2.32 & -1.76 & 0 & 0 \end{pmatrix} \begin{pmatrix} 0 \\ 0 \\ 300 \\ 300 \end{pmatrix} \\
&= -0.4\times1300 + 0.6\times 0 = -520
\end{align*}
et
\begin{align*}
E_1 &= 0.4\pi_1^T T(\xi_1) + 0.6\pi_2^T T(\xi_2) \\
& = 0.4  \begin{pmatrix} 0 & -3 & 0 -13 \end{pmatrix}
\begin{pmatrix} -60 & 0 \\ 0 & -80 \\ 0 & 0 \\ 0 & 0  \end{pmatrix} \\
& \quad + 0.6 \begin{pmatrix} -2.32 & -1.76 & 0 & 0 \end{pmatrix}
\begin{pmatrix} -60 & 0 \\ 0 & -80 \\ 0 & 0 \\ 0 & 0  \end{pmatrix} \\
& = -0.4  \begin{pmatrix} 0 & 240 \end{pmatrix}
-0.6 \begin{pmatrix} 139.2 & 140.8 \end{pmatrix} \\
& =  \begin{pmatrix} 0 & -96 \end{pmatrix} +
\begin{pmatrix} -83.52 & -84.48 \end{pmatrix} \\
& =  \begin{pmatrix} -83.52 & -180.48 \end{pmatrix}
\end{align*}
La coupe d'optimalité à générer est donc
$$
\theta \geq -520 - 83.52x_1 - 180.48x_2
$$
On retourne à la résolution du programme maître, qui est maintenant
\begin{align*}
\min\ & 100x_1 + 150x_2  + \theta \\
\mbox{t.q. }
& x_1 + x_2 \leq 120 \\
& x_1 \geq 40,\ x_2 \geq 20 \\
& \theta \geq -520 - 83.52x_1 - 180.48x_2
\end{align*}
\item
% http://opim.wharton.upenn.edu/~sok/papers/h/Higle-Tutorials2005-chapter02.pdf
(Tutorials in Operations Research, INFORMS, 2005)
La Dakota Furniture Company fabrique des bureaux, des tables et des chaises.
La fabrication de chaque type de mobilier nécessite du bois d’\oe{}uvre et deux types de main-d’œuvre qualifiée: la finition et la menuiserie. Le tableau ci-dessous indique le coût de chaque ressource et le montant nécessaire à la fabrication de chaque type de mobilier.

\begin{tabular}{|c|c|c|c|c|}
	\hline
	& Coût (\$) & Bureau & Table & Chaise \\
	\hline
	Bois de sciage (pieds-planche) & 2 & 8 & 6 & 1 \\
	\hline
	Finition (heures) & 4 & 4 & 2 & 1.5 \\
	\hline
	Menuiserie (heures) & 5.2 & 2 & 1.5 & 0.5 \\
	\hline
\end{tabular}

%Demand 150 125 300
Un bureau se vend à 60\$, une table à 40\$ et une chaise à 10\$.
Nous considérons trois scénarios de demande.

\begin{tabular}{|c|c|c|c|}
	\hline
	& Valeur basse & Valeur la plus probable & Valeur élevée \\
	\hline
	Bureaux & 50 & 150 & 250 \\
	\hline
	Tables & 20 & 110 & 250 \\
	\hline
	Chaises & 200 & 225 & 500 \\
	\hline
	Probabilité & 0.3 & 0.4 & 0.3 \\
	\hline
\end{tabular}

Afin de maximiser le profit total, quelle quantité de chaque article devrait être produite et quels sont les besoins en ressources correspondants?
Formulez le problème en tant que programme stochastique à deux étapes.
Remarque: seule la formulation est demandée. Vous n'avez pas à résoudre le problème.

Nous traitons le problème comme un problème d'allocation de ressources.

Variables:
\begin{itemize}
	\item $x_1$: quantité de bois de sciage
	\item $x_2$: nombre d'heures de finition
	\item $x_3$: nombre d'heures de menuiserie
	\item $y_1$: nombre de bureaux produits
	\item $y_2$: nombre de tables produites
	\item $y_3$: nombre de chaises produites
\end{itemize}

Formulation
\begin{align*}
\min_x \ & 2x_1 + 4x_2 + 5.2x_3 + \EE[Q(x,\bxi) ]\\
\mbox{t.q. } & x_1, x_2, x_3 \geq 0
\end{align*}
avec
\begin{align*}
Q(x,\xi) = \min_y \ & -60y_1-40y_2-10y_3 \\
\mbox{t.q.} \\
& 8y_1 + 6y_2 + y_3 \leq x_1 \\
& 4y_1 + 2y_2 + 1.5y_3 \leq x_2 \\
& 2y_1 + 1.5y_2 + 0.5y_3 \leq x_3 \\
& y_1 \leq d_1(\xi) \\
& y_2 \leq d_2(\xi) \\
& y_3 \leq d_3(\xi) \\
& y_1, y_2, y_3 \geq 0
\end{align*}

En considérant les trois scénarios de l'énoncé, nous avons
$$
\EE[Q(x,\bxi) ] = 0.3 Q(x,\xi_1)+ 0.4 Q(x,\xi_2) + 0.3 Q(x,\xi_3)
$$
où $\xi_1$, $\xi_2$ et $\xi_3$ correspondent aux valeurs basse, la plus probable et élevée, respectivement.
Dans ce cas, $d_1(\xi_1) = 50$, $d_2(\xi_1) = 20$, $d_3(\xi_1) = 200$, $d_1(\xi_2) = 150$, $d_2(\xi_2) = 110$, $d_3(\xi_2) = 225$, $d_1(\xi_3) = 250$, $d_2(\xi_3) = 250$, $d_3(\xi_3) = 500$.
\end{enumerate}

\chapter{Monte Carlo approximations}

\begin{enumerate}
\item
\begin{exercise}
Montrez que $E_{\bxi}[ \min_{x \in \cX} z(x,\bxi) ] \leq \min_{x \in \cX} E_{\bxi}[ z(x,\bxi) ]$, en considérant un ensemble réalisable $\cX$ déterministe.
Comment interpréter ce résultat par rapport à la programmation stochastique?
\end{exercise}

Considérons que les espérances sont bien définies et que les minimas existent, comme implicitement supposé dans l'énoncé.
Dès lors, il existe $x^* \in \cX$ tel que
$$
\min_{x \in \cX} E_{\bxi}[ z(x,\bxi) ] = E_{\bxi}[ z(x^*,\bxi) ].
$$
Comme $x^* \in \cX$,
$$
\min_{x \in \cX} z(x,\bxi) \leq z(x^*,\bxi)
$$
et dès lors
$$
E_{\bxi} \left[\min_{x \in \cX} z(x,\bxi)\right] \leq E_{\bxi} [z(x^*,\bxi)].
$$

%Remarquons que les ensembles réalisables étant identiques, s'il existe $x$ tel que $E_{\bxi} [z(x,\bxi)] < +\infty$, alors $E_{\bxi}[ \inf_{x \in \cX} z(x,\bxi) ] < +\infty$ et inversement, $E_{\bxi}[ \inf_{x \in \cX} z(x,\bxi) ] = +\infty$ implique $E_{\bxi} [z(x,\bxi)] = +\infty$.
Même si les minimas n'existent pas, l'énoncé reste valide en remplaçant les opérateurs de minimisation pas des infimums:
$$
E_{\bxi}\left[ \inf_{x \in \cX} z(x,\bxi) \right] \leq \inf_{x \in \cX} E_{\bxi}[ z(x,\bxi) ]
$$
En effet, 
$$
\inf_{x \in \cX} z(x,\bxi) \leq z(y,\bxi),\ \forall\, y \in \cX.
$$
Dès lors
$$
E_{\bxi}\left[ \inf_{x \in \cX} z(x,\bxi) \right] \leq E_{\bxi} [ z(y,\bxi) ],\ \forall\, y \in \cX,
$$
et donc l'énoncé est valide.

Le résultat indique que la minimisation sous information parfaite donne donne un estimateur biaisé de la solution.
Nous pouvons étendre le résultat à un échantillonnage Monte Carlo.
Soit
$$
z_N(x) = \frac{1}{N} \sum_{n = 1}^{N} z(x,\xi_n).
$$
Nous avons
\begin{align*}
E_{\bxi}\left[ \min_{x \in \cX} z_N(x) \right]
& = E_{\bxi}\left[ \min_{x \in \cX} \frac{1}{N} \sum_{n = 1}^{N} z(x,\bxi_n) \right] \\
& = \frac{1}{N} \sum_{n = 1}^{N} E_{\bxi}\left[ \min_{x \in \cX} z(x,\bxi_n) \right] \\
& \leq \frac{1}{N} \sum_{n = 1}^{N} \min_{x \in \cX} E_{\bxi}[ z(x,\bxi) ] \\
& = \min_{x \in \cX} E_{\bxi}[ z(x,\bxi) ]
\end{align*}

\end{enumerate}

\end{document}